\section{Introduction}
For a top-three prediction system, a pedigree representation is useful to the extent that it improves prediction beyond information already available from recent performance. Assessing that increment is difficult when training horses have direct graph vectors but many later entrants do not. A fallback based on embedded ancestors then changes the representation used for those entrants. Applying one construction rule across splits addresses that difference in procedure, although it does not make coverage or feature distributions identical.

We examine Japan Racing Association (JRA) flat races using a fixed training-root pedigree graph and chronological partitions \cite{jraResults,netkeibaRace,netkeibaDb,grover2016}. The central question is whether the specified graph-feature pipeline adds predictive value beyond race context and recent history when ancestor-based construction is used throughout, and how it compares with a simple parent-performance control. Four pipelines use the same cohort, three classifier families and separate turf and dirt models. The two primary comparisons are the ancestor-rule graph against history and against the parent control; the comparison with the original direct-else-ancestor pipeline is secondary. The test cohort had already been examined before this design was fixed, so these comparisons are post-test exploratory. The study evaluates particular feature pipelines rather than a new graph architecture or a causal effect of ancestry. We retain the earlier retrospective market analysis as separate context; the new four-arm study contains no market predictors.

\section{Related Work}
Sugiura \cite{sugiura2026} studied temporally held-out JRA win and place-rule outcomes using recent history, sire and damsire information and matched feature-block comparisons, without final odds or popularity. The place-rule label is not identical to the fixed top-three outcome here. A 2024 Japanese technical report proposes a pedigree graph neural network for future racehorse performance \cite{oikawa2024}; its detailed target and controls could not be established from the accessible abstract. Gulum \cite{gulum2018} used graph-derived features for paired-horse prediction. Oda and Onogi \cite{oda2024} studied pedigree-based breeding values for adjusted velocity and graded-winner status, different prediction targets from ours. Chung \emph{et al.} \cite{chung2024} considered race ranking in Seoul; we rank entry-classifier scores and do not evaluate a dedicated ranking model. Earlier racing studies also motivate history and market controls \cite{silverman2013,pudaruth2013}, while pedigree research motivates, but does not establish, incremental predictive value \cite{bailey2022}.

Our bounded contribution is the same-cohort evaluation of a fixed ancestor-rule graph pipeline against history and a specified chronological parent-outcome comparator, with the original mixed-representation pipeline as a secondary reference. Temporal validation, feature-block comparisons, race bootstrap and pedigree graphs each have prior precedents; none is claimed as a new method here.
